\section{Nine elliptic factors and rank constraints}\label{sec:nine-factors}

The four roots in \eqref{eq:common-four-roots} give nine elliptic
quotients.  Their rational twists can be identified explicitly, and the
images of a single strict point impose rank conditions on every factor.
Throughout this section put
\[
 s=t-1,\qquad T=t+1,\qquad d=sT,\qquad z=\ell+T.
\]
Thus the coordinate denoted by $\ell$ satisfies
\begin{equation}\label{eq:nine-roots}
 P^2=z,\qquad C^2=z-2t,\qquad
 M^2=-\frac{s(z-2T)}{z-T},\qquad
 B^2=\frac{T(z-2)}{z-T}.
\end{equation}
The smooth projective curve with this function field is $D_t$.
For a strict rational-distance point, its four roots satisfy
\[
 t>1,\qquad s<\ell<T,\qquad P,B>0,\qquad CM\ne0.
\]
Here $CM\ne0$ also follows directly from the square: $C=0$ would give
$x=y$, and $M=0$ would give $x+y=1$.  Either diagonal forces a positive
rational distance to be a nonzero rational multiple of $\sqrt2$.

\subsection{Character quotients and their rational models}

Define the elliptic curves
\begin{equation}\label{eq:nine-elliptic-models}
\begin{aligned}
 E_d &: Y^2=X(X-d)(X+d),\\
 G_t &: Y^2=X(X-d^2)(X+d),\\
 F_+ &: Y^2=X(X-2d)(X+2td),\\
 F_- &: Y^2=X(X-2d)(X-2td),\\
 H_+ &: Y^2=X(X-2T)(X-T^2),\\
 H_- &: Y^2=X(X-2s)(X-d).
\end{aligned}
\end{equation}
Each curve carries the point at infinity as its origin.  All six cubics
are nonsingular for $t\notin\{0,1,-1\}$.

\begin{theorem}\label{thm:nine-isogeny}
For every $t\in\Q\setminus\{0,1,-1\}$ there is an isogeny over $\Q$
\begin{equation}\label{eq:nine-isogeny}
 \Jac(D_t)\sim_{\Q}
 E_d^2\times G_t\times F_+^2\times F_-^2\times H_+\times H_-.
\end{equation}
The rank of the group of rational N\'eron--Severi classes of $\Jac(D_t)$
is at least twelve.
\end{theorem}
\begin{proof}
The sign changes of the four roots give a group
$V=(\Z/2\Z)^4$, in the order $(P,C,M,B)$.  On the $z$-line its six
inertia generators, over $0,2t,\infty,2T,2,T$, are
\[
 e_1,\quad e_2,\quad e_1+e_2,\quad
 e_3,\quad e_4,\quad e_3+e_4.
\]
A nonzero character supported on just one of the blocks
$\langle e_1,e_2\rangle$ and $\langle e_3,e_4\rangle$ is ramified at
two points.  Its double-cover quotient is rational.  There are six such
characters.  Each of the nine characters nonzero on both blocks is
ramified at four points and has a genus-one quotient.

The latter characters are the products of one member of
$(P,C,PC)$ and one member of $(M,B,MB)$.  Their elliptic models are
\begin{equation}\label{eq:nine-character-table}
\begin{array}{c|ccc}
 &M&B&MB\\\hline
 P&E_d&H_+&F_+\\
 C&H_-&E_d&F_-\\
 PC&F_-&F_+&G_t.
\end{array}
\end{equation}
Here are birational maps establishing all the models and their rational
twists.  The character in the first column specifies the genus-one
function field $\Q(z,\text{character})$:
\begin{equation}\label{eq:nine-map-table}
\begin{array}{c|c|c|c}
\text{character}&\text{target}&X&Y\\\hline
 PM&E_d&-s\ell&sPM\ell\\[1mm]
 BC&E_d&T\ell&TBC\ell\\[1mm]
 PB&H_+&Tz&TPB\ell\\[1mm]
 CM&H_-&s(2T-z)&sCM\ell\\[1mm]
 PCB&F_+&4tT\ell/z&4tT^2\ell PCB/z^2\\[1mm]
 CMB&F_-&4td/(T-\ell)&-4td\ell CMB/(T-\ell)^2\\[1mm]
 PCMB&G_t&dT(2t/z-1)&dtT\ell PCMB/z^2.
\end{array}
\end{equation}
For example, the $PB$ and $PCMB$ equations before transformation are
\[
 (PB\ell)^2=Tz(z-2)(z-T),\qquad
 (PCMB\ell)^2=-dz(z-2t)(z-2)(z-2T).
\]
Substitution of the corresponding rows gives the stated cubics.  The
other rows follow by multiplying their radicands in
\eqref{eq:nine-roots}.  Every displayed abscissa is an invertible
M\"obius function of $z$, and its ordinate is a nonzero rational function
of $z$ times the indicated character.  It therefore recovers that
character and defines a birational map of genus-one curves.  All maps
extend uniquely between the smooth projective models.  The inverse
image of the point at infinity is a rational point on the smooth
quotient; this point supplies the origin used for its Jacobian.
More explicitly, the origins in the seven rows lie above
$z=\infty,\infty,\infty,\infty,0,2T,0$, respectively.
Each is the unique ramified point of its genus-one double cover above
the indicated value.

The involution
\begin{equation}\label{eq:nine-iota}
 (\ell,P,C,M,B)\longmapsto(d/\ell,B,M,C,P)
\end{equation}
exchanges the $PCB$ and $PMB$ characters, and exchanges $CMB$ and
$PCM$.  Composing the two corresponding rows of
\eqref{eq:nine-map-table} with this involution gives
\[
\begin{aligned}
 PMB\longrightarrow F_+:\quad (X,Y)
 &=\left(\frac{4td}{\ell+s},
           \frac{4td\ell PMB}{(\ell+s)^2}\right),\\
 PCM\longrightarrow F_-:\quad (X,Y)
 &=\left(\frac{4ts\ell}{\ell-s},
          -\frac{4ts^2\ell PCM}{(\ell-s)^2}\right).
\end{aligned}
\]
Their origins lie above $z=2$ and $z=2t$, respectively.  This gives
the two remaining entries of \eqref{eq:nine-character-table}.

The pullback of the differential of each elliptic quotient lies in its
one-dimensional character eigenspace in $H^0(D_t,\Omega^1)$.
These nine distinct eigenspaces span that space, since $g(D_t)=9$.
The product norm map from $\Jac(D_t)$ to the nine quotient
Jacobians has the direct sum of these differential pullbacks as its
cotangent map.  It therefore has invertible differential and is an
isogeny over $\Q$.  Its dual gives the product pullback isogeny and
proves \eqref{eq:nine-isogeny}.

For an elliptic square $E^2$, the two factor divisors and the diagonal
have intersection matrix
\[
 \begin{pmatrix}0&1&1\\1&0&1\\1&1&0\end{pmatrix},
 \qquad\det=2.
\]
They give three independent rational divisor classes.  Apply this to
$E_d^2,F_+^2,F_-^2$, and take the polarization class on each of
$G_t,H_+,H_-$.  The twelve classes are independent on the product.
Pullback and the norm along a rational isogeny identify N\'eron--Severi
spaces after tensoring with $\Q$, proving the assertion.
\end{proof}

The genus-three quotient and its three Pryms also follow from the same
character decomposition.  Let $K\subset V$ be generated by the sign
changes of $(P,M)$ and of $(C,B)$.  None of the six inertia subgroups
meets $K$ nontrivially.  Thus $D_t\to D_t/K$ is an unramified cover
of degree four and its quotient has genus three.  Its invariant field
is generated by
\[
 u=\frac{PM-BC}{2},\qquad v=\frac{PM+BC}{2},
 \qquad z=\frac{u^2+v^2+2}{2}+t-\frac{uv}{t}.
\]
Consequently it is the smooth model of
\begin{equation}\label{eq:nine-quartic}
 C_t:\quad t^2(u^2+v^2)^2-4u^2v^2-4t^2d=0.
\end{equation}
The characters trivial on $K$ give
\[
 \Jac(C_t)\sim_{\Q}E_d^2\times G_t.
\]
The three intermediate double covers of $C_t$ have genus five.
Writing $\kappa_1=e_1+e_3$ and $\kappa_2=e_2+e_4$, their Pryms are
\[
\begin{aligned}
 \operatorname{Prym}(D_t/\langle\kappa_1\rangle\to C_t)
   &\sim_{\Q}F_+\times F_-,\\
 \operatorname{Prym}(D_t/\langle\kappa_2\rangle\to C_t)
   &\sim_{\Q}F_+\times F_-,\\
 \operatorname{Prym}(D_t/\langle\kappa_1+\kappa_2\rangle\to C_t)
   &\sim_{\Q}H_+\times H_-.
\end{aligned}
\]
The additional characters in these three quotients are respectively
$(PMB,PCM)$, $(PCB,CMB)$, and $(PB,CM)$.  Removing the three
characters of $C_t$ from each quotient proves the displayed isogenies.

\subsection{Two-descent and four auxiliary point sets}

We recall the elementary forms of two-descent used below; see also
\cite{cassels1991,silverman2009}.  On a split cubic
\[
 E:\quad y^2=(x-e_1)(x-e_2)(x-e_3)
\]
with distinct rational roots, the Kummer homomorphism has coordinates
$[x-e_i]$ in $\Q^\times/\Q^{\times2}$.  At $(e_i,0)$ its $i$th
coordinate is $[(e_i-e_j)(e_i-e_k)]$; at another root $(e_j,0)$ it is
$[e_j-e_i]$, and at infinity it is $1$.  The product of the three
coordinates is $1$.  The map has kernel $2E(\Q)$.
Indeed, the tangent-and-chord law proves multiplicativity, and when
all three coordinates are squares one may choose rational roots
$r_i^2=x-e_i$ with $r_1r_2r_3=-y$; the usual halving formulas give a
point with abscissa
\[
 x+r_1r_2+r_1r_3+r_2r_3.
\]
The same calculation at a root shows that $(e_i,0)$ has a rational half
if and only if both $e_i-e_j$ and $e_i-e_k$ are squares.

For the integral model $y^2=x^3+ax^2+bx$, the one-coordinate
isogeny descent sends an ordinary point to $[x]$, sends $(0,0)$ to
$[b]$, and sends infinity to $1$.  Its kernel is the image of the dual
isogeny from
\[
 E':\quad y^2=x^3-2ax^2+(a^2-4b)x.
\]
Its image consists of signed squarefree divisors $q$ of $b$ for which
\begin{equation}\label{eq:nine-isogeny-quartic}
 W^2=qU^4+aU^2V^2+(b/q)V^4,\qquad\gcd(U,V)=1,
\end{equation}
is soluble in integers, together with the stated values at zero and
infinity.  To obtain this equation write $x=q(U/V)^2$ and substitute
in the cubic.  A prime not dividing $b$ cannot occur to an odd exponent
in $x$: negative valuations follow from the leading cubic term, and
positive valuations from the constant term of $x^2+ax+b$.
These observations give the support assertion and the homogeneous
quartic.

If $E$ has full rational two-torsion, write $\alpha(E)$ and
$\alpha'(E')$ for these two images.  Then
\begin{equation}\label{eq:nine-rank-descent}
 2^{\rank E(\Q)}=\frac{|\alpha(E)|\,|\alpha'(E')|}{4}.
\end{equation}
For completeness, the dual compositions are multiplication by two.
The other rational two-torsion points of $E$ map onto the kernel of
the dual isogeny.  The index calculation therefore gives
$[E(\Q):2E(\Q)]=|\alpha(E)|\,|\alpha'(E')|$.
Mordell--Weil and $E(\Q)[2]\simeq(\Z/2\Z)^2$ identify the left side
with $2^{\rank E(\Q)+2}$, giving
\eqref{eq:nine-rank-descent}.

\begin{lemma}\label{lem:nine-small-curves}
The following are complete rational point sets:
\begin{align*}
 E_1(\Q)&=\{O,(0,0),(1,0),(-1,0)\},\\
 E_2(\Q)&=\{O,(0,0),(2,0),(-2,0)\},\\
 A(\Q)&=\{O,(0,0),(6,0),(8,0)\},
   &&A:\ y^2=x(x-6)(x-8),\\
 E_0(\Q)&=\{O,(0,0),(-1,0),(-4,0),(2,\pm6),(-2,\pm2)\},
   &&E_0:\ y^2=x(x+1)(x+4).
\end{align*}
Here $E_1$ and $E_2$ are the curves $y^2=x^3-x$ and $y^2=x^3-4x$.
\end{lemma}
\begin{proof}
For $E_1$, the first descent image is $\{1,-1\}$.
Its dual $y^2=x^3+4x$ has nonnegative rational abscissae, so its image
is contained in $\{1,2\}$.  Both classes occur, the latter at $(2,4)$.
Equation~\eqref{eq:nine-rank-descent} gives rank zero.
For $E_2$, the first image is $\{1,-1,2,-2\}$, all represented by
two-torsion.  Its dual $y^2=x^3+16x$ has only the candidates $1,2$.
The class $2$ would require
\[
 W^2=2U^4+8V^4.
\]
If $U$ is odd the right side has valuation one at $2$; otherwise $V$
is odd and it has valuation three.  Thus the second image is
$\{1\}$ and the rank is zero.

For every positive rational $d$, the torsion of $E_d$ is its rational
two-torsion.  A rational half of the largest root $d$ would require
both $d$ and $2d$ square; the other roots have negative differences.
The third-division polynomial
\[
 3x^4-6d^2x^2-d^4
\]
cannot vanish at a rational number, since it would give
$x^2/d^2=1\pm2\sqrt3/3$.  Mazur's torsion theorem
\cite{mazur1977}, in the full-two-torsion case, leaves only
$(\Z/2\Z)^2$ after orders four and three have been excluded.
This completes the point lists for $E_1,E_2$.

For $A$, the model has $a=-14,b=48$.  There are no real points with
$x<0$, so the first descent image is contained in $\{1,2,3,6\}$;
the four two-torsion points realize all four classes.  The dual is
$y^2=x^3+28x^2+4x$, with candidates $1,-1,2,-2$.
For $q=-1$, equation~\eqref{eq:nine-isogeny-quartic} is nonsquare
modulo three for every primitive pair.  For $q=2$ its right side is
\[
 2\bigl((U^2-V^2)^2+16U^2V^2\bigr);
\]
opposite parity gives valuation one at two, and both odd gives
valuation five.  For $q=-2$ the right side is
\[
 -2\bigl((U^2-V^2)^2-12U^2V^2\bigr),
\]
with valuations one and three, respectively.  Thus the dual image
is $\{1\}$, proving rank zero.

For $E_0$, the first descent image is $\{1,-1,2,-2\}$, as witnessed
by $O,(-1,0),(2,6),(-2,2)$.  Its dual is
$y^2=x(x-1)(x-9)$ and has only the candidates $1,3$.
The second class would give
\[
 W^2=3U^4-10U^2V^2+3V^4.
\]
If both variables are units at three the right side is $2$ modulo
three; if just one is a unit its valuation is one.  Thus the dual
image is $\{1\}$ and the rank is zero.

To determine the last two torsion groups, the following table gives
all cubic values at the good primes five and seven, in the order
$x=0,1,\ldots,p-1$:
\[
\begin{array}{c|c|c|c}
\text{curve}&p&f(0),\ldots,f(p-1)&\#E(\F_p)\\\hline
 A&5&0,0,3,0,2&4\\
 A&7&0,0,6,3,4,1,0&8\\
 E_0&5&0,0,1,4,0&8\\
 E_0&7&0,3,1,0,6,4,0&8.
\end{array}
\]
The nonzero squares are $1,4$ modulo five and $1,2,4$ modulo seven.
Each point count follows by adding one for infinity and counting the
two square roots of each nonzero square.  Prime-to-$p$ torsion injects
under good reduction.  Combining the two primes bounds the torsion
orders by four and eight, respectively: the five-primary part is
excluded using seven, and the seven-primary part using five.
The displayed rational points attain these bounds, completing both
point lists.
\end{proof}

\subsection{Torsion of the six factors}

We shall need a convenient parametrization of a rational flex.
\begin{lemma}\label{lem:nine-flex}
If $y^2=x(x-A)(x-B)$ is nonsingular and has a rational point of order
three, then there is $k\in\Q\setminus\{0,1,-1,-\tfrac12,-2\}$ such that
\begin{equation}\label{eq:nine-flex}
 \frac BA=\frac{1+2k}{(1-k)(1+k)^3}.
\end{equation}
There is also a nonzero rational $r$ with $A=r^2(1-k^2)$.
\end{lemma}
\begin{proof}
Write the tangent at the order-three point $(x,y)$ as
$h(X)=m(X-x)+y$.  Since it meets the cubic with multiplicity three,
\[
 X(X-A)(X-B)-h(X)^2=(X-x)^3.
\]
For each root $e\in\{0,A,B\}$, the rational number
$r_e=h(e)/(x-e)$ satisfies $r_e^2=x-e$.  These numbers are nonzero
and distinct.  The identity $h(e)=y-mr_e^2=r_e^3$ shows that they are
the three roots of $R^3+mR^2-y$.  Hence their pairwise products sum
to zero.  Put $r=r_0$ and $k=r_A/r$.  Then
$r_B/r=-k/(1+k)$, and
\[
 A=r^2(1-k^2),\qquad B=r^2\frac{1+2k}{(1+k)^2}.
\]
This gives \eqref{eq:nine-flex}.  The values $0,\pm1$ contradict
nonzero roots or $A\ne0$; $-1/2$ makes $B=0$, and $-2$ makes $A=B$.
\end{proof}

\begin{lemma}\label{lem:nine-three-torsion}
For rational $t\notin\{0,1,-1\}$, the curve $G_t$ has no rational
three-torsion.  For $t>1$ the following additional assertions hold:
\begin{enumerate}
\item if $T$ is a rational square, neither $H_+$ nor $F_-$ has
rational three-torsion;
\item if $T$ is twice a rational square, $F_+$ has no rational
three-torsion.
\end{enumerate}
\end{lemma}
\begin{proof}
Set
\[
 D=(1-k)(1+k)^3,\quad R=\frac{1+2k}{D},\quad
 N=2+4k-2k^3-k^4.
\]
Then
\begin{equation}\label{eq:nine-flex-shifts}
 R-1=\frac{k^3(k+2)}D,\qquad R+1=\frac ND.
\end{equation}
Suppose first that $H_+$ has rational three-torsion and $T$ is a
square.  Take $A=2T,B=T^2$ in Lemma~\ref{lem:nine-flex}.
It gives $T/2=R$ and
\[
 r^2=\frac{4(1+2k)}{(1-k)^2(1+k)^4}.
\]
Thus $1+2k=a^2$ and $1-k^2=2b^2$ for rational $a,b$, whence
$8b^2=3+2a^2-a^4$.  Write $a=U/V$ in lowest terms.  The rational
number $2bV^2$ is an integer, since its square multiplied by two is
integral.  Calling it $W$ gives
\[
 2W^2=3V^4+2U^2V^2-U^4.
\]
Modulo sixteen the right side is $4$ if $U,V$ are both odd, is
$7$ or $15$ if only $U$ is odd, and is $3$ or $11$ if only $V$
is odd.  None is $0,2,8$, the possible residues of $2W^2$.

For $F_-$ use $A=2d,B=2td$, so $t=R$.  If $T$ is square,
\eqref{eq:nine-flex-shifts} and $A=r^2(1-k^2)$ show that
\[
 \frac{r^2}T=
 \left(\frac{k}{(1-k)(1+k)^2}\right)^2 2k(k+2).
\]
Hence $2k(k+2)$ is square.  If $v_2(k)<0$ its valuation is
$1+2v_2(k)$, and if $v_2(k)=0$ its valuation is one.  Consequently
$v_2(k)>0$.  But then $D$ is a dyadic unit and $N\equiv2\pmod8$;
thus $v_2(T)=v_2(N/D)=1$, contradicting that $T$ is square.

For $F_+$ take $A=2d,B=-2td$, so $-t=R$.  If $T$ is twice square,
\eqref{eq:nine-flex-shifts} gives
\[
 -2k(k+2)(1-k^2)
 =2T\left(\frac{(1-k)(1+k)^2}{k}\right)^2\in\Q^{\times2}.
\]
On the quartic $w^2=-2k(k+2)(1-k^2)$, the transformation
\[
 X=4-4/k,\qquad Y=-4w/k^2
\]
gives the curve $A$ of Lemma~\ref{lem:nine-small-curves}.
The branch point $k=0$ maps to infinity.  The quartic has no rational
points at infinity, since its leading coefficient is $2$.
The complete point list on $A$ therefore leaves just
$k=0,-2,1,-1$, all excluded by Lemma~\ref{lem:nine-flex}.

Finally, rational three-torsion on $G_t$ would give
$A=d^2,B=-d$, and hence $R=-1/d$.  Since $d+1=t^2$,
\[
 t^2=\frac{R-1}R=\frac{k^3(k+2)}{1+2k}.
\]
Consequently
\[
 (X,Y)=\left(2k,\frac{2t(1+2k)}k\right)
\]
is a point on $E_0$.  Its complete list in
Lemma~\ref{lem:nine-small-curves} leaves
$k\in\{0,-\tfrac12,-2,1,-1\}$, again impossible.
\end{proof}

\subsection{Positive rank on every factor}

\begin{theorem}\label{thm:nine-positive-ranks}
Let a strict rational-distance point have product parameter
$t=r_0r_3+r_1r_2$, and let the curves be as in
\eqref{eq:nine-elliptic-models}.  Each of
$E_d,G_t,F_+,F_-,H_+,H_-$ has positive rational rank.  In particular,
\begin{equation}\label{eq:nine-rank-nine}
 \rank\Jac(D_t)(\Q)\ge9.
\end{equation}
\end{theorem}
\begin{proof}
Every ordinate in \eqref{eq:nine-map-table} is nonzero at the strict
point.  The point $(-s\ell,sPM\ell)$ on $E_d$ therefore has infinite
order, by the torsion argument in
Lemma~\ref{lem:nine-small-curves}.  If $d$ were square, $E_d$ would
be rationally isomorphic to $E_1$; if $2d$ were square it would be
isomorphic to $E_2$.  Both are rank zero.  Thus
\begin{equation}\label{eq:nine-d-exclusions}
 d\notin\Q^{\times2},\qquad 2d\notin\Q^{\times2}.
\end{equation}

We next exclude rational order four on the remaining curves.
A two-torsion point can have a rational half only at the largest
root, because any other root has a negative root difference.
For $F_+$ its half would require $2d$ square.  For $H_+$ and $H_-$
it would require $d$ square, after removing square factors.
For $G_t$ it would require $d^2+d=dt^2$ square.  These contradict
\eqref{eq:nine-d-exclusions}.

For $F_-$ the largest root is $2td$.  Its half requires both
$2td$ and $2d(t-1)$ square.  Their quotient gives
$t/(t-1)$ square, so $b^2=t(t-1)$ for a nonzero rational $b$;
the second condition also gives $a^2=2(t+1)$.
Then $(X,Y)=(a^2,2ab)$ lies on
\[
 Y^2=X(X-2)(X-4)
\]
with nonzero ordinate.  Translation $X\mapsto X-2$ identifies this
curve with $E_2$, contradicting its point list.

Mazur's theorem now implies that torsion on each of the five curves
is either $(\Z/2\Z)^2$ or $\Z/2\Z\times\Z/6\Z$.
Lemma~\ref{lem:nine-three-torsion} excludes the second case for $G_t$,
so its nonzero-ordinate image is of infinite order.
For the other four curves, odd torsion is divisible by two.  Hence its
Kummer image is trivial, and the Kummer image of all torsion is just
the image of the rational two-torsion.

The common roots yield the identities
\begin{equation}\label{eq:nine-point-kummer}
\begin{aligned}
 X_{F_+}-2d&=2(TC/P)^2,&
 X_{F_-}-2d&=2(sB/M)^2,\\
 X_{H_+}&=TP^2,&
 X_{H_-}-2s&=-sC^2.
\end{aligned}
\end{equation}
On $H_-$, the possible torsion classes of the coordinate $X-2s$
are $1$ and $[-2s]$.  The actual class $[-s]$ differs from the first
by sign and from the second by $[2]$.  This image has infinite order.
On $H_+$, the torsion classes of $X$ are $1,[2T]$.
The actual class $[T]$ differs from both unless $T$ is square.
In that case Lemma~\ref{lem:nine-three-torsion} excludes order three,
and the nonzero ordinate excludes two-torsion.

For $F_-$, the torsion classes of $X-2d$ are
\[
 1,\quad[-2d],\quad[-s],\quad[2T].
\]
The actual class $[2]$ can meet this list only when $T$ is square;
the same lemma handles that case.  For $F_+$ the list is
\[
 1,\quad[-2d],\quad[T],\quad[-2dT].
\]
The actual class $[2]$ meets it only when $T$ is twice square,
which is the remaining case of that lemma.  Thus all six ranks are
positive.  Their multiplicities in \eqref{eq:nine-isogeny} give
$2+1+2+2+1+1=9$.
\end{proof}

\subsection{Two independent images and the rank strata}

Two quotient maps can furnish independent points on a single elliptic
curve.  The following argument uses the full Kummer rows of the points
and of their sum.  We write pairs of rational numbers for their classes
in $(\Q^\times/\Q^{\times2})^2$.

On $E_d$, take the two diagonal images
\[
 Q=(-s\ell,sPM\ell),\qquad R=(T\ell,TBC\ell).
\]
For the coordinates $(X,X-d)$ their rows are
\begin{equation}\label{eq:nine-cm-rows}
 \delta Q=(-s\ell,-s),\qquad
 \delta R=(T\ell,T),\qquad
 \delta(Q+R)=(-d,-d).
\end{equation}
For example $X(Q)-d=-sP^2$ and $X(R)-d=TC^2$.
The four torsion rows are
\begin{equation}\label{eq:nine-cm-torsion}
 (1,1),\quad(-1,-d),\quad(d,2),\quad(-d,-2d).
\end{equation}
Since $d>0$ is nonsquare, the sum row in
\eqref{eq:nine-cm-rows} is absent from this list.
Comparing the individual rows, with signs retained, shows that either
one can meet a torsion row only if $T$ is square or twice square.
Thus
\begin{equation}\label{eq:nine-cm-two}
 T\notin\Q^{\times2}\cup2\Q^{\times2}
 \quad\Longrightarrow\quad \rank E_d(\Q)\ge2.
\end{equation}
Indeed the images of $Q,R$ in
$E_d(\Q)/(2E_d(\Q)+E_d(\Q)_{\mathrm{tors}})$ are then nonzero and
distinct, hence linearly independent over $\F_2$.

For $F_+$ and $F_-$, take the image in
\eqref{eq:nine-map-table} and its image after the involution
\eqref{eq:nine-iota}.  For both pairs the coordinates $(X,X-2d)$ give
\begin{equation}\label{eq:nine-f-rows}
 \delta Q_\pm=(tT\ell,2),\qquad
 \delta R_\pm=(ts\ell,2),\qquad
 \delta(Q_\pm+R_\pm)=(d,1).
\end{equation}
The first-coordinate products are, explicitly,
\[
 X(Q_+)X(R_+)=d\left(\frac{4tT}{PB}\right)^2,\qquad
 X(Q_-)X(R_-)=d\left(\frac{4ts}{MC}\right)^2.
\]
The second coordinates follow from \eqref{eq:nine-point-kummer}.
The complete torsion rows are
\begin{equation}\label{eq:nine-f-torsion}
\begin{array}{c|cccc}
 F_+&(1,1)&(-t,-2d)&(2d,T)&(-2td,-2dT)\\
 F_-&(1,1)&(t,-2d)&(2d,-s)&(2td,2T).
\end{array}
\end{equation}
For $F_+$, the sum row $(d,1)$ never occurs: the positive rows
would require $d$ or $2$ to be square.  An individual row in
\eqref{eq:nine-f-rows} can occur only if $T$ is twice square.
It follows that
\begin{equation}\label{eq:nine-fplus-two}
 T\notin2\Q^{\times2}\quad\Longrightarrow\quad\rank F_+(\Q)\ge2.
\end{equation}
For $F_-$, an individual row can meet torsion only when $T$ is
square.  The sum row can meet torsion only if both $2t$ and $2T$
are squares.  Therefore
\begin{equation}\label{eq:nine-fminus-two}
\begin{gathered}
 T\notin\Q^{\times2},\qquad
 \text{at least one of $2t,2T$ is nonsquare}
 \\
 \Longrightarrow\quad\rank F_-(\Q)\ge2.
\end{gathered}
\end{equation}

\begin{theorem}[Rank strata]\label{thm:nine-rank-strata}
For the product parameter of any strict rational-distance point, the
following lower bounds hold:
\begin{equation}\label{eq:nine-rank-table}
\begin{array}{l|ccc|c}
 \text{condition}&\rank E_d&\rank F_+&\rank F_-&\rank\Jac(D_t)\\\hline
 T\notin\Q^{\times2}\cup2\Q^{\times2}&2&2&2&15\\
 T\in\Q^{\times2}&1&2&1&11\\
 T\in2\Q^{\times2},\ 2t\notin\Q^{\times2}&1&1&2&11\\
 T\in2\Q^{\times2},\ 2t\in\Q^{\times2}&1&1&1&9.
\end{array}
\end{equation}
Every entry in the table is a lower bound.
\end{theorem}
\begin{proof}
Combine \eqref{eq:nine-cm-two}, \eqref{eq:nine-fplus-two}, and
\eqref{eq:nine-fminus-two} with
Theorem~\ref{thm:nine-positive-ranks}.  Formula~\eqref{eq:nine-isogeny}
gives the rank identity
\[
 2\rank E_d+\rank G_t+2\rank F_++2\rank F_-+
 \rank H_++\rank H_-.
\]
Substitution gives the final column.  The four rows exhaust the
positive rational values of $T$, since $2$ is nonsquare.
\end{proof}

The theorem gives a useful exclusion criterion for a fixed parameter:
a rank bound below the relevant entry in
\eqref{eq:nine-rank-table} excludes a strict point on that fibre.
It also places every fibre arising from a strict point outside the
classical Chabauty inequality $\rank\Jac(D_t)(\Q)<g(D_t)$.
